\section{Coding product families and observable tests}\label{sec:C23:catalog}
Table~\ref{tab:C23:catalog} places coding, terminal, editor, and application-building systems into one comparison program without treating them as equivalent. A product may span local and remote execution, and an exposed client does not make its entire service inspectable. The official surfaces were rechecked on 29 October 2026; experimental claims require the corresponding executable configuration and evidence. Retrieval failure on a product homepage is not a capability verdict: Lovable's official documentation establishes its application-building and code-export surfaces, while the resolved Antigravity 2.0 announcement distinguishes a standalone agent application from its separate IDE~\cite{cumb2026lovable,cumb2026antigravity2}.
\begingroup\setlength{\tabcolsep}{3pt}\small
\bookneedspace{.42\textheight}
\begin{longtable}{L{25mm}L{35mm}L{48mm}}
\caption{Coding and computer-work product surfaces and diagnostic tests.}\label{tab:C23:catalog}\\
\toprule Product & Inspected role & Diagnostic condition \\
\midrule\endfirsthead
\toprule Product & Inspected role & Diagnostic condition \\
\midrule\endhead
Codex\cite{executionproduct005} & Open harness and SDK boundaries & Bind final artifacts to tested revisions; separate local execution from hosted service state. \\
OpenHands\cite{executionproduct008} & Self-hosted software-agent workflows & Inspect sandbox reset and replay after environment loss. \\
goose\cite{executionproduct009} & Local extensible tool automation & Audit extension authority and effects outside the workspace. \\
Warp\cite{executionproduct013} & Terminal and cloud development surfaces & Compare local state with remotely delegated job state at interruption. \\
Kilo\cite{executionproduct015} & Multi-model coding interfaces & Hold the selected model fixed before attributing gains to orchestration. \\
OpenCode\cite{executionproduct016} & Extensible coding runtime & Change tool schemas and test stored session compatibility. \\
Letta Code / Letta harness\cite{executionproduct017} & Stateful coding and editable memory & Revoke an old instruction and test its influence on resumed work. \\
Claude Code\cite{executionproduct021} & Terminal coding and developer integration & Verify permission enforcement independently of model-generated intent. \\
GitHub Copilot\cite{executionproduct022} & Repository editor and delegated coding surfaces & Distinguish IDE suggestions from autonomous remote task execution. \\
Cline\cite{executionproduct023} & Configurable local coding automation & Inject a failing command and inspect confirmation and retry behavior. \\
Zed Agent\cite{executionproduct025} & Editor with native and external agents & Test unsaved buffers, external-agent context and artifact identity. \\
Cursor\cite{executionproduct027} & Editor and cloud coding interfaces & Audit which project state reaches background workers and their final merge. \\
Sourcegraph Agentic Batch Changes\cite{executionproduct028} & Batch repository-change workflow & Measure migration coverage and per-repository regressions rather than one patch score. \\
Hermes Agent\cite{executionproduct030} & Persistent agent with memory and skills & Separate stored procedures from learned weight updates and test stale-skill effects. \\
Mistral Vibe\cite{executionproduct036} & Coding CLI and work interfaces & Identify local command mediation and managed conversation boundaries. \\
Devin\cite{executionproduct037} & Delegated engineering environments & Verify reproducible environment setup and independently test the delivered patch. \\
Factory Droids\cite{executionproduct038} & Local and cloud engineering work & Track workspace reuse, credential scope and review artifact provenance. \\
Amp\cite{executionproduct050} & Persistent coding threads & Interrupt work after tool dispatch and inspect repeated effects on resumption. \\
Antigravity\cite{cumb2026antigravity2} & Standalone agent app and separate IDE & Test project-scoped context, subagents and permissions under the selected product surface. \\
Qwen Code\cite{cumb2026qwencode} & Coding runtime with selectable model adapters & Match model and tool serialization before comparing repository results. \\
mini-swe-agent\cite{executionproduct065} & Minimal research coding loop & Use as a low-complexity baseline with identical evaluator and budget. \\
Junie\cite{executionproduct067} & IDE, CLI and delegated coding workflow & Separate test execution evidence from generated explanations of correctness. \\
Harness Cosmos / Auggie\cite{executionproduct069} & Enterprise code-context tooling & Provider lineage does not prove shared implementation; test repository-context boundaries. \\
Kiro\cite{executionproduct070} & Specification-oriented development & Perturb a requirement after planning and test invalidation of dependent work. \\
Rovo Dev\cite{executionproduct073} & Issue-context engineering workflow & Verify that ticket permissions and repository write permissions remain distinct. \\
Replit Agent\cite{executionproduct074} & Application construction and deployment & Validate deployed behavior and environment configuration, not preview appearance alone. \\
MiniMax Code\cite{executionproduct075} & Persistent mini-app and task surfaces & Test background-task lifecycle and whether saved artifacts match verified output. \\
Tabnine Agent / Plugin for OpenCode\cite{executionproduct086,cumb2026tabnineopencode} & Enterprise coding distribution & Verify selected binary, provider lockdown and tool-governance configuration. \\
v0\cite{executionproduct087} & Generated full-stack web applications & Test backend access controls, deployment identity and post-generation regressions. \\
Lovable\cite{cumb2026lovable} & Full-stack app building and code export & Bind preview, exported code and published deployment to the same reviewed artifact. \\
CodeBuddy\cite{executionproduct095} & Subagent interface documentation & Test isolated contexts, inherited permissions and integration of conflicting edits. \\
TRAE / SOLO\cite{executionproduct096} & IDE and standalone task agent & Separate interactive assistance from unattended completion under the same task contract. \\
ZCode\cite{executionproduct097} & Desktop coding with tool interfaces & Audit memory persistence and browser authority separately from coding quality. \\
\bottomrule\end{longtable}
\endgroup


The important mechanisms cut across these families. Editor integration can reduce navigation cost and expose project context, but it can also hide unsaved state or privileged extensions. Cloud delegation can provide isolated environments and parallel work, but adds environment reconstruction and artifact-transfer boundaries. Application builders can integrate deployment tightly, but task success must then include the deployed artifact and its operational configuration rather than only a generated preview.

Codex's documented SDK and open harness interfaces provide concrete inspection boundaries, while its hosted application and services remain separate components~\cite{executionproduct005}. Claude Code's public repository and SDK documentation likewise require artifact-specific interpretation~\cite{executionproduct021,executionproduct020}. No claim of complete implementation visibility follows from a public issue tracker or client repository. For both families, experiments must identify what runs locally, what runs remotely, and where state and authority are enforced.

The current inventory also contains consequential context-lifecycle distinctions. Letta Code's pinned README describes memory, skills, prompts and harness modifications as adaptation surfaces; its cloud backend manages bundled default prompts on resume unless the documented preservation option is selected. That is a configuration transition requiring capture, not evidence of a fixed policy or autonomous weight improvement~\cite{cumb2026lettacode}. Qwen Code exposes multiple provider protocols and SDK/desktop/browser interfaces, so treating every Qwen Code run as one Qwen model is incorrect~\cite{cumb2026qwencode}. Tabnine's OpenCode distribution supplies its own managed binary, Tabnine-served model selection and organization controls; sharing an underlying agent family does not establish equivalence to a separately configured OpenCode installation~\cite{cumb2026tabnineopencode}. Provider lineage and product packaging must remain distinct from the manifest of an evaluated run.

In the pinned Codex TypeScript SDK, \texttt{CodexExec.run} spawns the CLI with a JSON event stream, and \texttt{Thread.run} reduces completed-item events into a buffered return value. A turn-failure event causes that buffered path to throw; completed items accumulated earlier are not returned as a successful \texttt{Turn}. The streaming path exposes events as they arrive~\cite{cumb2026codexsdk}. An application needing partial-progress recovery should durably record relevant events and artifact identities while streaming, rather than infer that a thrown turn produced no work. Neither stream persistence nor a completed-item event by itself certifies the absence of unresolved remote effects. This is a source-level distinction between event delivery and task completion, not an independently measured failure rate.

